\documentclass[11pt]{article}
\ifdefined\pdfmapfile\pdfmapfile{+lm.map}\pdfmapfile{+cm.map}\pdfmapfile{+symbols.map}\fi
\usepackage[T1]{fontenc}
\usepackage{lmodern}
\usepackage[margin=1in]{geometry}
\usepackage{amsmath,amssymb,amsthm,hyperref}
\hypersetup{colorlinks=true,linkcolor=blue,urlcolor=blue,citecolor=blue,pdftitle={No finite positive Renyi order has a universal completely monotone heat-entropy derivative},pdfauthor={Keio Research Notes}}
\setlength{\emergencystretch}{2em}
\newtheorem{theorem}{Theorem}
\newtheorem{lemma}[theorem]{Lemma}
\newtheorem{corollary}[theorem]{Corollary}
\newcommand{\av}[1]{\left\langle #1\right\rangle}
\newcommand{\R}{\mathbb R}
\title{No finite positive R\'enyi order has a universal completely monotone heat-entropy derivative}
\author{Research manuscript prepared for Keio}
\date{15 November 2026}
\begin{document}
\maketitle

\paragraph{Scope and status.}
This note provides an elementary spectral obstruction and a real-line localization proof. External review is pending. It does not claim that the Shannon case is new: that case was disproved by Gu--Sellke (2026). The periodic/Gaussian-window method also has prior precedent in Zou--Fan--Gao--Wang (2026); an anonymous Evidence Press candidate (6 November 2026) applies a related method to classify Tsallis orders. The proposed additional statement here is the pointwise classification for \emph{R\'enyi} entropy: there are no surviving finite positive orders. Priority has not been established by an exhaustive literature review. No log-concavity assertion is made, and the input density and failing derivative may depend on the entropy order.

Wu--Yu--Guo's Question 1 asks which R\'enyi orders satisfy the universal all-derivative property. The theorem below classifies the admissible finite positive orders for that property: the set is empty. It does not settle their Question 2, which asks for the admissible entropy-order range for each fixed derivative order.

Put $G_s(x)=(2\pi s)^{-1/2}\exp(-x^2/(2s))$ and $P_t f=G_t*f$. For $\alpha>0$, $\alpha\ne1$, put
\[
 h_\alpha(f)=\frac{1}{1-\alpha}\log\int_\R f^\alpha,
 \qquad h_1(f)=-\int_\R f\log f.
\]

\begin{theorem}\label{thm:main}
For every fixed finite $\alpha>0$ there are $s>0$, an integer $m\ge1$, and $t_0>0$, such that the probability density in the fixed family
\[
 f_s(x)=\frac{G_s(x)(1+\tfrac14\cos x+\tfrac1{16}\cos2x)}{Z_s},
 \qquad Z_s=1+\tfrac14 e^{-s/2}+\tfrac1{16}e^{-2s},
\]
satisfies
\[
 (-1)^{m-1}\left.\frac{d^m}{dt^m}h_\alpha(P_t f_s)\right|_{t=t_0}<0.
\]
In particular no fixed finite positive R\'enyi order has a completely monotone first entropy derivative along every one-dimensional heat flow. The counterexamples are smooth, strictly positive, and bounded above and below by positive multiples of a Gaussian.
\end{theorem}

\paragraph{1. A spectral obstruction.}
Write $\av v=(2\pi)^{-1}\int_0^{2\pi}v(\theta)\,d\theta$. The periodic heat flow
\[
 u_t(\theta)=1+\varepsilon e^{-t/2}\cos\theta+\delta e^{-2t}\cos2\theta
\]
has mean one. Let $H_\alpha(t)=(1-\alpha)^{-1}\log\av{u_t^\alpha}$ when $\alpha\ne1$, with $H_1(t)=-\av{u_t\log u_t}$.

\begin{lemma}\label{lem:spectral}
Suppose $F(t)=\sum_{j\ge1}c_j e^{-jt}$ with $\sum_{j\ge1}j|c_j|<\infty$ and with termwise differentiation valid on $t>0$. If $c_r>0$ for some integer $r$, then for all sufficiently large integers $k$,
\[
 (-1)^k F^{(k+1)}(k/r)<0.
\]
In particular $F'$ is not completely monotone.
\end{lemma}
\begin{proof}
Termwise differentiation gives
\[
 \frac{(-1)^kF^{(k+1)}(k/r)}{r^k e^{-k}}
 =\sum_{j\ge1}(-jc_j)\left[\frac jr\exp\left(1-\frac jr\right)\right]^k.
\]
The bracket lies in $(0,1]$, and equals $1$ only when $j=r$, since $\log x\le x-1$ with equality only at $x=1$. Dominated convergence, using $\sum j|c_j|<\infty$, shows that the displayed expression converges to $-rc_r<0$.
\end{proof}

Fix any $\varepsilon,\delta\ge0$ with $\varepsilon+\delta<1$ and put $z=e^{-t/2}$. The function
\[
 M_\alpha(z)=\av{(1+\varepsilon z\cos\theta+\delta z^4\cos2\theta)^\alpha}
\]
is analytic on a sufficiently small disk $|z|<R_\alpha$, and is uniformly close to $1$ on a smaller closed disk about zero. Indeed the expression inside the power is uniformly close to $1$, so the power and logarithm use their analytic branches at $1$. Its invariance under $(z,\theta)\mapsto(-z,\theta+\pi)$ implies an even power series. Therefore, for all sufficiently large $T=T_\alpha>0$,
\[
 H_\alpha(T+t)=\sum_{j\ge1}(c_j e^{-jT})e^{-jt},\qquad \sum_{j\ge1}j|c_j|e^{-jT}<\infty.
\]
The same statements hold directly for $H_1=-\av{u\log u}$. The coefficients at $\alpha=1$ also equal the limits of the coefficients at other orders, because $M_1(z)=1$ identically and the quotient defining $H_\alpha$ has a removable singularity.

Writing $w=z^2$, the first three coefficients of $M_\alpha$ are
\begin{align*}
 [w]M_\alpha&=\frac{\alpha(\alpha-1)}4\varepsilon^2=:a_1,\\
 [w^2]M_\alpha&=\frac{\alpha(\alpha-1)(\alpha-2)(\alpha-3)}{64}\varepsilon^4=:a_2,\\
 [w^3]M_\alpha&=\frac{\alpha(\alpha-1)(\alpha-2)}8\varepsilon^2\delta
 +\frac{\alpha(\alpha-1)(\alpha-2)(\alpha-3)(\alpha-4)(\alpha-5)}{2304}\varepsilon^6=:a_3.
\end{align*}
These identities use $\av{\cos^2\theta}=1/2$, $\av{\cos^4\theta}=3/8$, $\av{\cos^6\theta}=5/16$, and $\av{\cos^2\theta\cos2\theta}=1/4$. No other monomial has the needed spectral weight and a nonzero average.

Since the first three coefficients of $\log M_\alpha$ are $a_1$, $a_2-a_1^2/2$, and $a_3-a_1a_2+a_1^3/3$, respectively, we obtain the exact identities
\begin{align}
 c_2&=\frac{\alpha(\alpha^2+3\alpha-6)}{64}\varepsilon^4,\label{eq:c2}\\
 c_3&=\frac{\alpha(2-\alpha)}8\varepsilon^2\delta
 -\frac{\alpha(\alpha^4+4\alpha^3-4\alpha^2-25\alpha+30)}{576}\varepsilon^6.\label{eq:c3}
\end{align}
They remain valid at $\alpha=1$ by the removable extension.

Now fix once and for all $\varepsilon=1/4$ and $\delta=1/16=\varepsilon^2$. If $\alpha\ge3/2$, then $\alpha^2+3\alpha-6\ge3/4$, so $c_2>0$. If $0<\alpha\le3/2$, write $Q(\alpha)=\alpha^4+4\alpha^3-4\alpha^2-25\alpha+30$. The elementary bound $|Q(\alpha)|\le144$ on $[0,2]$ gives
\[
 c_3=\alpha\varepsilon^4\left(\frac{2-\alpha}{8}-\frac{Q(\alpha)\varepsilon^2}{576}\right)
 \ge\alpha\varepsilon^4\left(\frac1{16}-\frac1{64}\right)>0.
\]
Thus the same fixed periodic density works at every positive finite order. The time shift $T_\alpha$ multiplies $c_j$ by the positive number $e^{-jT_\alpha}$ and preserves its sign. Lemma~\ref{lem:spectral} applied to $F(t)=H_\alpha(T_\alpha+t)$ supplies a strict sign violation at some finite $m$ and positive $t_0$ in every case.

\paragraph{2. Gaussian localization.}
The following lemma transfers this finite-order strict violation to $\R$.

\begin{lemma}\label{lem:local}
For this lemma allow any fixed $\varepsilon,\delta\ge0$ with $\varepsilon+\delta<1$, and define $f_s=G_s(1+\varepsilon\cos x+\delta\cos2x)/Z_s$ with $Z_s=1+\varepsilon e^{-s/2}+\delta e^{-2s}$. For every fixed finite $\alpha>0$, every integer $m\ge1$, and every compact interval $I\subset(0,\infty)$,
\[
 \frac{d^m}{dt^m}h_\alpha(P_t f_s)\longrightarrow H_\alpha^{(m)}(t)
\]
uniformly on $I$ as $s\to\infty$.
\end{lemma}
\begin{proof}
Gaussian convolution gives, for every integer $k$,
\[
 P_t(G_s e^{ik\cdot})(x)=G_{s+t}(x)
 \exp\left(-\frac{stk^2}{2(s+t)}\right)
 \exp\left(\frac{iksx}{s+t}\right).
\]
Put $v=s+t$, $\omega=s/v$, $A_s=\varepsilon\exp(-st/(2v))$, $B_s=\delta\exp(-2st/v)$, and $U_s(\theta,t)=1+A_s(t)\cos\theta+B_s(t)\cos2\theta$. Then
\[
 P_t f_s(x)=Z_s^{-1}G_v(x)U_s(\omega x,t).
\]
For $\alpha\ne1$, Gaussian integration shows
\[
 \int(P_t f_s)^\alpha
 =Z_s^{-\alpha}(2\pi v)^{(1-\alpha)/2}\alpha^{-1/2}Q_{s,\alpha}(t),
\]
where
\[
 Q_{s,\alpha}(t)=\mathbb E[U_s(\Theta_s,t)^\alpha],
 \qquad \Theta_s\sim N\left(0,\frac{s^2}{\alpha v}\right).
\]
Consequently
\begin{equation}\label{eq:renyi-real}
 h_\alpha(P_t f_s)=\frac12\log(2\pi v)+\frac{\log\alpha}{2(\alpha-1)}
 -\frac{\alpha}{1-\alpha}\log Z_s+\frac{\log Q_{s,\alpha}(t)}{1-\alpha}.
\end{equation}

Here are details justifying convergence of every required time derivative. There is a fixed strip $|\Im\theta|\le\rho$, $\rho>0$, on which $|A_s\cos\theta+B_s\cos2\theta|\le q<1$, uniformly for $t\in I$ and sufficiently large $s$. All fixed finite time derivatives of $A_s,B_s$ are uniformly bounded and converge to those of $\varepsilon e^{-t/2},\delta e^{-2t}$, respectively. Thus $U_s^\alpha$, and its first $m$ time derivatives, are uniformly analytic and bounded on this strip. If $b_{s,k}(t)$ are their Fourier coefficients before differentiation, the contour-shift estimate gives
\[
 |\partial_t^j b_{s,k}(t)|\le C_j e^{-\rho|k|},\qquad 0\le j\le m.
\]
Fourier integration now yields the absolutely convergent series
\[
 Q_{s,\alpha}(t)=\sum_{k\in\mathbb Z}b_{s,k}(t)
 \exp\left(-\frac{k^2s^2}{2\alpha(s+t)}\right).
\]
For $s$ large relative to $I$, the exponent is at most $-c_\alpha k^2s$. Each fixed time derivative of its exponential is bounded by $C_m(1+k^2)^m\exp(-c_\alpha k^2s)$: this follows from the chain rule and the bounded derivatives of $s^2/(s+t)$ of order at least one. The nonzero-frequency sum and its first $m$ derivatives therefore tend uniformly to zero. The zero-frequency coefficient tends in $C^m(I)$ to $\av{u_t^\alpha}$, which is bounded away from zero. Therefore $\log Q_{s,\alpha}\to\log\av{u_t^\alpha}$ in $C^m(I)$. For every positive derivative order, the derivative of $\frac12\log v$ tends to zero, while $Z_s$ is constant in $t$. Equation~\eqref{eq:renyi-real} proves the assertion for $\alpha\ne1$.

For $\alpha=1$, write $p=P_t f_s$. Directly,
\[
 h_1(p)=\frac12\log(2\pi v)+\frac{1}{2v}\int x^2p(x)\,dx
 +\log Z_s-Z_s^{-1}\mathbb E_{X\sim N(0,v)}[U_s(\omega X,t)\log U_s(\omega X,t)].
\]
The same Fourier argument, now for $U_s\log U_s$, shows that the final expectation tends in $C^m(I)$ to $\av{u_t\log u_t}$. The Gaussian second-moment identity
\[
 \mathbb E_{X\sim N(0,v)}[X^2\cos(k\omega X)]
 =v(1-k^2\omega^2v)\exp(-k^2\omega^2v/2)
\]
shows that $v^{-1}\int x^2p(x)dx\to1$ in $C^m(I)$: the nonconstant terms, and all their fixed time derivatives, are polynomially bounded factors times an exponential tending to zero. Finally $Z_s\to1$ and every positive derivative of $\frac12\log v$ tends to zero. This proves the Shannon assertion as well.
\end{proof}

Apply Lemma~\ref{lem:local} at the finite $m,t_0$ supplied by the circle argument. The strict negative signed derivative persists for all sufficiently large finite $s$, proving Theorem~\ref{thm:main}.

\paragraph{An explicitly located periodic example.}
At $\alpha=2$, with $\delta=0$, one has $H_2(t)=-\log(1+ce^{-t})$, where $c=\varepsilon^2/2$. Setting $u=ce^{-t}$ gives
\[
 -H_2^{(4)}(t)=\frac{u(1-4u+u^2)}{(1+u)^4}.
\]
Thus $\varepsilon=9/10$ and $t=1/10$ give $-H_2^{(4)}(t)<0$, since $u=(81/200)e^{-1/10}>2-\sqrt3$. The localization proof gives a real-line fourth-derivative violation for every sufficiently large $s$. An explicit finite $s$ can also be certified using the exact formula
\[
 h_2(P_t f_s)=\frac12\log(s+t)-\log F_s(t)+C_s,
\]
\[
 F_s(t)=1+2\varepsilon\exp\left(-\frac{st}{2(s+t)}-\frac{s^2}{4(s+t)}\right)
 +\frac{\varepsilon^2}{2}\exp\left(-\frac{st}{s+t}\right)
 \left[1+\exp\left(-\frac{s^2}{s+t}\right)\right],
\]
where $C_s$ is independent of $t$. For the explicit finite choice $s=1000$, the accompanying script \texttt{verify\_renyi.py} evaluates the fourth Taylor coefficient in Arb ball arithmetic, at both 96-bit and 160-bit precision. Both computations certify the rational enclosure
\[
 \frac{337}{10000}<\left.\frac{d^4}{dt^4}h_2(P_t f_{1000})\right|_{t=1/10}<\frac{338}{10000}.
\]
Here this finite example uses $\varepsilon=9/10$, $\delta=0$, hence
\[
 f_{1000}(x)=\frac{e^{-x^2/2000}(1+\tfrac9{10}\cos x)}{\sqrt{2000\pi}(1+\tfrac9{10}e^{-500})}.
\]
The approximate derivative is $0.0337233593130162$. The script verifies this finite computation and the displayed coefficient identities, not the general analytic proof. The universal-order theorem above does not depend on the numerical certificate.

\paragraph{Relation to the corpus methods.}
The proof uses exact low-order algebra, spectral sign isolation, exponential domination, and controlled asymptotic localization. The signed spectral coefficient is the obstruction; localization is needed to make it a counterexample to the real-line probabilistic question. These tools are used where the proof needs them; this note does not claim that all ten source-corpus proofs are necessary ingredients.

\begin{thebibliography}{9}
\bibitem{WYG} H. Wu, L. Yu, L. Guo, \emph{On the Complete Monotonicity of R\'enyi Entropy}, IEEE Transactions on Information Theory 71(7), 4895--4914 (2025), \href{https://arxiv.org/abs/2312.01819}{arXiv:2312.01819}. The universal question appears as Question 1 in v3.
\bibitem{GS} Y. Gu, M. Sellke, \emph{A Counterexample to the Gaussian Completely Monotone Conjecture}, \href{https://arxiv.org/abs/2605.11656}{arXiv:2605.11656} (2026).
\bibitem{ZFGW} Zou, Fan, Gao, Wang, \emph{An Explicit Family of Log-Concave Counterexamples to the Gaussian Completely Monotone Conjecture}, \href{https://arxiv.org/abs/2608.30275}{arXiv:2608.30275} (2026).
\bibitem{EP} Anonymous, \emph{Complete monotonicity of Tsallis entropy under heat flow: a sharp order classification}, Evidence Press, unrefereed candidate v0.1.0 (6 November 2026), \href{https://doi.org/10.5281/zenodo.24001602}{doi:10.5281/zenodo.24001602}. This is a different entropy functional; its related classification and methods must not be claimed as new here.
\end{thebibliography}
\end{document}
